\section*{Capítulo \ref{ch:g11:ocean-ridges} --- As dorsais oceânicas e a litosfera oceânica}

\begin{solution}{exo:g11:ocean-ridges:1}
O \omterm{def:g11:ocean-ridges:peridotite}{peridotito} é uma rocha escura, de grão grosso, feita principalmente de
olivina, com um pouco de piroxênio; grande parte do manto é feita dele. O
\omterm{def:g11:ocean-ridges:gabbro}{gabro} é uma rocha ígnea escura, de grão grosso, de feldspato plagioclásio e
piroxênio, a rocha mais abundante da \omterm{def:g11:ocean-ridges:oceanic-lithosphere}{crosta oceânica} profunda. O \omterm{def:g11:ocean-ridges:gabbro}{gabro}
pertence à crosta, o \omterm{def:g11:ocean-ridges:peridotite}{peridotito} ao manto. O \omterm{def:g11:ocean-ridges:gabbro}{gabro} cristalizou lentamente, no
fundo, sob o fundo do mar, onde o magma esfriou devagar: os seus cristais
tiveram tempo para crescer. A lava basáltica esfriou depressa no fundo do
mar, e os seus cristais ficaram pequenos.
\end{solution}

\begin{solution}{exo:g11:ocean-ridges:2}
Do fundo do mar para baixo: sedimentos (nenhum no eixo), lavas basálticas
(almofadas e derrames em lençol), complexo de diques, \omterm{def:g11:ocean-ridges:gabbro}{gabro} e, abaixo do
\omterm{def:g11:seismic-structure:moho}{Moho}, o \omterm{def:g11:ocean-ridges:peridotite}{peridotito} do manto. O furo atravessou $1507/7100 \approx 0.21$,
cerca de um quinto da crosta: os \omterm{def:g11:ocean-ridges:gabbro}{gabros} profundos e o \omterm{def:g11:seismic-structure:moho}{Moho} nunca foram
alcançados por uma sonda.
\end{solution}

\begin{solution}{exo:g11:ocean-ridges:3}
$d(0) = \qty{2500}{m}$; $d(9) = 2500 + 350 \times 3 = \qty{3550}{m}$;
$d(25) = 2500 + 350 \times 5 = \qty{4250}{m}$.
\end{solution}

\begin{solution}{exo:g11:ocean-ridges:4}
Uma \omterm{def:g11:ocean-ridges:normal-fault}{falha normal} é uma falha inclinada ao longo da qual o bloco acima do
plano de falha desceu em relação ao bloco abaixo. Numa dorsal lenta, as
placas estiram a \omterm{def:g11:lithosphere-asthenosphere:lithosphere}{litosfera} fria e frágil: ela é esticada, e se rompe. Os
blocos acima das falhas deslizam para baixo, em direção ao eixo, e o vale
axial é o bloco que desceu entre elas.
\end{solution}

\begin{solution}{exo:g11:ocean-ridges:5}
Ela não ferve por causa da pressão: mais de dois quilômetros de água pesam
sobre ela, e sob essa pressão a água continua líquida muito acima de
\qty{100}{\celsius}. A ``fumaça'' negra é uma nuvem de grãos minúsculos de
sulfetos de ferro e de outros minerais, que precipitam quando a água quente
da fonte, rica em metais dissolvidos, encontra a água fria do mar.
\end{solution}

\begin{solution}{exo:g11:ocean-ridges:6}
(1) A cerca de \qty{54}{km} de profundidade, a cerca de \qty{1324}{\celsius}.
(2) A cerca de \qty{85}{km} de profundidade, a cerca de $1400 + 0.45 \times 85 \approx
\qty{1438}{\celsius}$. (3) O manto mais quente começa a fundir mais fundo, de
modo que uma coluna mais longa de rocha que sobe funde antes de chegar à
crosta: ele produz mais basalto, e uma crosta mais espessa.
\end{solution}

\begin{solution}{exo:g11:ocean-ridges:7}
(1) $1489 - 1345 = \qty{144}{\celsius}$ abaixo da sua temperatura de fusão.
(2) Ela sobe $100 - 54 = \qty{46}{km}$; a diferença se fecha em
$144/46 \approx \qty{3.1}{\celsius}$ por quilômetro, em média. (3) Não: na
subida, a rocha até esfria um pouco, cerca de $0.45 \times 46 \approx
\qty{21}{\celsius}$. É a sua temperatura de fusão que cai, muito mais
depressa, à medida que a pressão diminui.
\end{solution}

\begin{solution}{exo:g11:ocean-ridges:8}
As semitaxas são de \qty{7.3}{cm} e \qty{1.65}{cm} por ano, isto é,
\qty{73}{km} e \qty{16.5}{km} por milhão de anos. Para produzir
\qty{1000}{km} de fundo de um lado: $1000/73 \approx 13.7$ milhões de anos
para a Dorsal do Pacífico Oriental, $1000/16.5 \approx 61$ milhões de anos
para a Dorsal Mesoatlântica.
\end{solution}

\begin{solution}{exo:g11:ocean-ridges:9}
(1) \qty{5}{cm} equivalem a \qty{5e-5}{km}: $65000 \times 5 \times 10^{-5} =
\qty{3.25}{km^2}$ de novo fundo por ano. (2) $3.25 \times 7.1 \approx
\qty{23}{km^3}$ de crosta por ano. (3) $3.6 \times 10^8 / 3.25 \approx 1.1
\times 10^8$ anos, cerca de 110 milhões de anos: nesse ritmo, as dorsais
poderiam produzir todo o fundo oceânico em cerca de cem milhões de anos, um
tempo curto em comparação com a idade da Terra. É por isso que nenhum fundo
oceânico é muito antigo.
\end{solution}

\begin{solution}{exo:g11:ocean-ridges:10}
(1) O plano de falha faz \ang{60} com a horizontal, e o rebaixamento é o
lado vertical de um triângulo retângulo cuja hipotenusa é o deslocamento:
$1.5/\sin\ang{60} \approx \qty{1.73}{km}$. (2) Cada falha afasta os lados
horizontalmente $1.5/\tan\ang{60} \approx \qty{0.87}{km}$; as duas falhas
juntas, cerca de \qty{1.73}{km}. (3) Não. Os \qty{25}{km} são a distância
entre as falhas que limitam o bloco rebaixado, não o estiramento, que é de
apenas cerca de \qty{1.7}{km}; e a dorsal continua a se abrir por novos
diques e lavas no eixo, além das falhas.
\end{solution}

\begin{solution}{exo:g11:ocean-ridges:11}
$d(25) = 6400 - 3200\,e^{-25/62.8} \approx 6400 - 3200 \times 0.672 \approx
\qty{4251}{m}$; $d(100) = 6400 - 3200\,e^{-100/62.8} \approx 6400 - 3200
\times 0.203 \approx \qty{5749}{m}$. Entre 25 e 100 milhões de anos, o
fundo afunda cerca de \qty{1500}{m}. A lei da raiz quadrada dá
$2500 + 350 \times 5 = \qty{4250}{m}$ a 25 milhões de anos: as duas leis
concordam onde ambas se aplicam.
\end{solution}

\begin{solution}{exo:g11:ocean-ridges:12}
(1) $\alpha\,\Delta T\,L = 3.2 \times 10^{-5} \times 700 \times 100 =
\qty{2.24}{km}$. (2) $s = 3380 \times 2.24/(3380 - 1025) = 7571/2355 \approx
\qty{3.2}{km}$. (3) A lei exponencial permite $6400 - 2500 =
\qty{3900}{m}$ de afundamento ao todo: o resfriamento e a contração sozinhos
explicam cerca de quatro quintos disso, com valores simples para a queda de
temperatura e a espessura.
\end{solution}

\begin{solution}{exo:g11:ocean-ridges:13}
(1) $54 - 7 = \qty{47}{km}$ contra $85 - 7 = \qty{78}{km}$: a coluna mais
quente é $78/47 \approx 1.66$ vez mais longa. (2) $7.1 \times 1.66 \approx
\qty{11.8}{km}$. (3) Os \qty{10.3}{km} medidos são mais espessos que o
normal, como previsto, mas menos que a estimativa: os quilômetros mais
profundos da coluna, logo acima do início da fusão, dão menos material
fundido que os mais rasos, de modo que a hipótese da mesma quantidade de
material fundido por quilômetro superestima.
\end{solution}

\begin{solution}{exo:g11:ocean-ridges:14}
(1) A lei dá \qty{3178}{m}, \qty{4160}{m}, \qty{4782}{m} e
\qty{5267}{m}; as medianas medidas são mais rasas em \qty{64}{m},
\qty{142}{m}, \qty{317}{m} e \qty{619}{m}. (2) A diferença cresce com a
idade. (3) A lei dá a profundidade do próprio fundo, corrigida dos seus
sedimentos; as células medidas carregam os seus sedimentos, que são mais
espessos no fundo mais antigo, e incluem montes submarinos e planaltos
vulcânicos, que são mais rasos que o fundo em volta.
\end{solution}

\begin{solution}{exo:g11:ocean-ridges:15}
(1) $p = 1025 \times 9.8 \times 2500 \approx \qty{25.1}{MPa}$, mais
\qty{0.1}{MPa}: cerca de \qty{25.2}{MPa}. (2) Isso fica acima da pressão
crítica, \qty{22.064}{MPa}: a água da fonte não pode ferver, por mais quente
que esteja. (3) A pressão chega a \qty{22.064}{MPa} à profundidade
$(22.064 - 0.1) \times 10^6/(1025 \times 9.8) \approx \qty{2190}{m}$: toda
fonte mais funda que cerca de \qty{2200}{m} está livre da ebulição.
\end{solution}

\begin{solution}{pb:g11:ocean-ridges:1}
\textbf{1.} $14.6/2 = \qty{7.3}{cm}$ por ano, isto é, \qty{73}{km} por
milhão de anos.

\textbf{2.} À profundidade do eixo da dorsal, $d(0) = \qty{2500}{m}$.

\textbf{3.} Do fundo do mar para baixo: lavas basálticas (almofadas e
derrames em lençol), complexo de diques e \omterm{def:g11:ocean-ridges:gabbro}{gabro}, cerca de \qty{7.1}{km} de
crosta ao todo; abaixo do \omterm{def:g11:seismic-structure:moho}{Moho}, o \omterm{def:g11:ocean-ridges:peridotite}{peridotito} do manto, até a base da
\omterm{def:g11:lithosphere-asthenosphere:lithosphere}{litosfera}, que ainda é muito fina no eixo. Os sedimentos vão se depositar
sobre ele mais tarde.

\textbf{4.} Quando as placas se afastam, a \omterm{def:g11:lithosphere-asthenosphere:lithosphere}{astenosfera} sobe sob a dorsal. A
sua temperatura cai apenas cerca de \qty{0.45}{\celsius} por quilômetro,
enquanto a sua temperatura de fusão cai muito mais depressa à medida que a
pressão diminui. A cerca de \qty{54}{km} de profundidade, as duas se
encontram: a rocha começa a fundir, e continua fundindo enquanto sobe, sem
que nenhum calor seja acrescentado (\omterm{def:g11:ocean-ridges:partial-melting}{fusão por descompressão}).

\textbf{5.} A cada milhão de anos, a dorsal produz $2 \times 73 = \qty{146}{km}$
de fundo dos seus dois lados; para cada quilômetro de dorsal, o volume de
crosta é $146 \times 7.1 \times 1 \approx \qty{1040}{km^3}$.

\textbf{6.} $73 \times t$: \qty{73}{km}, \qty{292}{km}, \qty{1168}{km} e
\qty{4672}{km}.

\textbf{7.} $2000/73 \approx 27$ milhões de anos.

\textbf{8.} As placas mudam de velocidade e de direção ao longo de dezenas de
milhões de anos, como mostra a curva da cadeia Havaí--Imperador
(\cref{ch:g11:plate-motions}); a taxa de hoje não precisa valer por 64
milhões de anos.

\textbf{9.} Longe do eixo, a crosta já não fica sobre uma lente de magma: ela
esfria, e já não há calor suficiente abaixo para impulsionar a circulação e
aquecer a água.

\textbf{10.} A oeste da Dorsal do Pacífico Oriental, o pedaço viaja na placa
do Pacífico. Para oeste, a placa do Pacífico termina ao longo de fossas
profundas, onde afunda sob as suas vizinhas: zonas de subducção
(\cref{ch:g11:subduction-zones}).

\textbf{11.} $d(1) = \qty{2850}{m}$, $d(4) = 2500 + 700 = \qty{3200}{m}$,
$d(16) = 2500 + 1400 = \qty{3900}{m}$, $d(64) = 2500 + 2800 = \qty{5300}{m}$.

\textbf{12.} $3900 - 2500 = \qty{1400}{m}$ durante os primeiros 16 milhões
de anos, e $5300 - 3900 = \qty{1400}{m}$ entre 16 e 64 milhões de anos: o
mesmo afundamento, num tempo três vezes mais longo. O fundo afunda cada vez
mais devagar à medida que envelhece.

\textbf{13.} $350\sqrt{t} = 1500$ dá $\sqrt{t} \approx 4.29$, logo
$t \approx 18.4$ milhões de anos.

\textbf{14.} $6400 - 3200\,e^{-64/62.8} \approx 6400 - 3200 \times 0.361
\approx \qty{5245}{m}$, cerca de \qty{55}{m} a menos que a lei da raiz
quadrada: as duas leis concordam com diferença de cerca de \qty{1}{\percent}.

\textbf{15.} A mediana medida é cerca de \qty{650}{m} mais rasa que a lei. A
lei descreve o fundo corrigido do peso dos seus sedimentos; as
profundidades medidas incluem os sedimentos, muitas dezenas de milhões de
anos deles num fundo assim, e os montes submarinos e os planaltos
vulcânicos, e tudo isso torna o fundo medido mais raso.

\textbf{16.} $\alpha\,\Delta T = 3.2 \times 10^{-5} \times 700 = 0.0224$: a
parte mantélica fica \qty{2.24}{\percent} mais densa,
$3380 \times 1.0224 \approx \qty{3456}{kg/m^3}$.

\textbf{17.} A parte mantélica tem $100 - 7.1 = \qty{92.9}{km}$ de espessura. Densidade
média: $(7.1 \times 2900 + 92.9 \times 3456)/100 \approx (20590 + 321060)/100
\approx \qty{3416}{kg/m^3}$.

\textbf{18.} $3416 > 3380$: a \omterm{def:g11:lithosphere-asthenosphere:lithosphere}{litosfera} antiga como um todo é cerca de
\qty{36}{kg/m^3}, cerca de \qty{1}{\percent}, mais densa que a \omterm{def:g11:lithosphere-asthenosphere:lithosphere}{astenosfera}
abaixo dela. A crosta leve já não basta para compensar a sua parte
mantélica fria e densa.

\textbf{19.} Sendo mais densa que a \omterm{def:g11:lithosphere-asthenosphere:lithosphere}{astenosfera}, ela pode afundar de volta
no manto numa fossa: a subducção, assunto do próximo capítulo.

\textbf{20.} Aos 64 milhões de anos, o pedaço de fundo fica a cerca de
\qty{5300}{m}: cerca de \textbf{5.3~km} de profundidade, \qty{2.8}{km} mais
fundo do que onde nasceu.
\end{solution}
